\chapter{Known parameter classes}

\section{Classical classes}

The classical route combines hyperbolic or subhyperbolic dynamics, puzzle
pieces, and the phase--parameter correspondence.  The main classes to record
are:
\begin{itemize}
\item \textbf{Interior parameters.}  Local connectivity is local in an
      open subset of the plane.
\item \textbf{Hyperbolic and subhyperbolic cases.}  Expansion or
      attracting-cycle control gives finite geometric models and parameter
      control; the foundational source is Douady--Hubbard.
\item \textbf{Misiurewicz and related postcritically finite cases.}
      Finite orbit combinatorics and landing data, treated in the
      Douady--Hubbard/Yoccoz framework.
\item \textbf{Finitely renormalizable cases.}  Divergent puzzle moduli and
      shrinking parameter pieces, with Yoccoz as the primary source.
\end{itemize}

The exact hypotheses matter.  The review will not use the informal phrase
``all regular parameters'' as a Lean theorem name; each class receives a
predicate and a proof obligation.

\section{Lyubich's machinery}

Lyubich's work supplies the renormalization and geometric language needed
after the finitely renormalizable case: principal nests, modulus estimates,
rigidity, a priori bounds, and parapuzzles.  The 1997 Acta Mathematica
papers are the core reference for the renormalization framework.  The
Astérisque paper
\emph{Dynamics of quadratic polynomials, III: parapuzzle and SBR measures}
is kept as a separate ``Volume III'' source in the inventory.

The phrase ``Volume IV'' is not used as a canonical identifier in the
current source inventory.  Later results are therefore cited by their actual
paper or preprint title until a primary bibliographic source establishes a
formal continuation of the series.

\section{Dudko--Lyubich results}

Two recent parameter classes are especially important for this project:
\begin{enumerate}
\item bounded-type satellite parameters, treated in the GAFA paper and
      \lean{arXiv:1808.10425};
\item Feigenbaum points, treated in arXiv:2309.02107 and the 2026
      Publications Mathématiques de l'IHÉS publication.
\end{enumerate}
Both results use substantial renormalization, a priori-bound, and
parameter-puzzle machinery.  The formalization target is not merely the
final sentence ``MLC holds at $c$'', but the chain of definitions and lemmas
that produces the shrinking connected parameter basis.

\section{The 2025--2026 expansion}

The current source inventory adds several results that belong in the
renormalization layer rather than in a generic bibliography:
\begin{center}\small
\begin{tabular}{@{}P{0.31\textwidth}@{\hspace{4pt}}P{0.48\textwidth}@{\hspace{4pt}}P{0.15\textwidth}@{}}
\textbf{Source} & \textbf{Formalization-relevant content} & \textbf{Role}\\
\hline
Kahn--Kapiamba--Lyubich,
\lean{arXiv:2606.27272} &
MLC and a priori bounds for parabolically bounded primitive
renormalization; Thin--Thick Decomposition, Value Calculus, Wanderers
Theorem, and Wave Lemma. & MLC case\\
Dudko--Lim,
\lean{arXiv:2609.15500} &
Uniform bounds for bubbles of neutral quadratic polynomials. & geometry\\
Dudko--Lim--Lyubich,
\lean{arXiv:2609.37383} &
Rigidity of the attractor of neutral renormalization and combinatorial
rigidity of neutral towers. & rigidity\\
Shen, \lean{arXiv:2610.25983} &
Complex analytic structure of hybrid classes of generalized
polynomial-like mappings. & analytic base\\
\end{tabular}
\end{center}

The first of these is a new pointwise MLC class and must eventually be
represented by a precise Lean predicate for parabolically bounded primitive
renormalization.  The neutral-renormalization papers do not by themselves
prove global MLC, but they supply geometric and rigidity lemmas that the
future parameter-puzzle and straightening layers must consume.  The hybrid
class theorem belongs one level below renormalization: it describes the
analytic parameter space in which polynomial-like straightening lives.

\section{Dudko's survey and the remaining frontier}

Dudko's arXiv:2512.24171 is used as a current status and roadmap source.
It is not treated as a replacement for the original proofs.  In particular,
the review separates:
\[
\text{survey statement}
\ne
\text{precise theorem}
\ne
\text{Lean proof}.
\]
